% HD-F1-T38. GPT-6 Astra Pro, /root/transport; 22 November 2026 UTC.
\section[HD-F1-T38: an effective odd symplectic rank cutoff]{HD-F1-T38: an effective natural-rank cutoff for every remaining odd-characteristic symplectic source}
\label{sec:f1-symplectic-bounded-level-rank-cutoff}

Put \(L_0=3\cdot10^{30}-3\), as in T36. Every original
remaining T2-OTHER source with almost simple projective
socle \(\operatorname{PSp}_{2r}(q)\), for odd \(q\), satisfies
\[
 \boxed{\qquad r\le L_0^2+4<10^{61}.\qquad}
 \tag{F1T38.1}
\]
The retained field bound is \(q\le27\). This gives a finite
range of natural ranks for this entire symplectic branch;
it does not assert the existence of a source in that range.
The argument retains every original field and diagonal
automorphism, scalar phase and restriction multiplicity.
The only numerical tests below supplement uniform proofs.

\paragraph{1. The original layer and a literal rank-two commutator.}
Retain the original finite group \(G\subset GL_K(V)\),
\(m=\dim_KV\ge27\), scalar centre \(Z\), perfect layer
\(L\subset G\), and the actual surjection
\(\beta:H=\operatorname{Sp}_{2r}(q)\longrightarrow L\).
We argue by contradiction with \(r\ge L_0^2+5\).
T36 and T37 apply in this range. In particular, every
irreducible constituent of \(V_{\mathbf C}|_H\) has
minimal tensor level \(k\) satisfying
\[
                         2\le k\le L_0.
 \tag{F1T38.2}
\]
Here the tensor generator is the actual direct sum of the
two reducible Weil modules. T37's exact four-label action
also shows that the presence of any level-one constituent
would put the entire original restriction in its excluded
level-one orbit.

For each original \(g\in G\), the actual lifted action is
\[
 \alpha_g(h)=a h^{p^e}a^{-1},\qquad
 a\in\operatorname{GSp}_{2r}(q),\qquad
 \beta\alpha_g=c_g\beta,
 \tag{F1T38.3}
\]
with the multiplier of \(a\) denoted by \(\lambda\).
These are precisely the central-cover and automorphism maps
proved in T37, using the finite-field Steinberg theorem;
the present rank excludes all cover exceptions. If
\(g\notin Z\), then \(\alpha_g\ne1\): otherwise its
projective image centralizes the simple socle in its
literal almost simple group, whose centralizer is trivial.

Such an \(\alpha_g\) moves a projective line of the natural
symplectic space \(A=\mathbf F_q^{2r}\). Indeed, a
semilinear map \(a\sigma_e\) fixing every line is diagonal
in the coordinate basis. The lines \(\langle e_i+e_j\rangle\)
make all its diagonal entries equal, and the lines
\(\langle e_i+t e_j\rangle\) make \(t^{p^e}=t\) for
every \(t\in\mathbf F_q\). It is therefore a scalar linear
map and induces the identity automorphism of \(H\).

Choose a nonzero \(v\) on a moved line, set
\(w=a(v^{p^e})\), and use the actual transvection
\(t_{v,c}(x)=x+c\langle x,v\rangle v\), with \(c\ne0\).
Direct substitution gives
\[
 \alpha_g(t_{v,c})=t_{w,c'},\qquad
 c'=c^{p^e}/\lambda\ne0,
 \qquad h=\alpha_g(t_{v,c})t_{v,c}^{-1}.
 \tag{F1T38.4}
\]
Writing \(\ell_u(x)=\langle x,u\rangle\), one has
\[
 (h-I)x=-c\ell_v(x)v+
 c'\bigl(\ell_w(x)-c\ell_v(x)\ell_w(v)\bigr)w.
 \tag{F1T38.5}
\]
The vectors \(v,w\) and the functionals \(\ell_v,\ell_w\)
are independent, and the displayed coefficient matrix
has determinant \(-cc'\ne0\). Consequently
\[
 \operatorname{rank}_A(h-I)=2,\qquad
 \ker(h-I)=v^\perp\cap w^\perp,
 \qquad
 \beta(h)=[g,\beta(t_{v,c})].
 \tag{F1T38.6}
\]
No coefficient-space inertia operator has been rescaled.

\paragraph{2. The actual tensor module and its code embeddings.}
Fix any one irreducible constituent \(W\) of minimal level
\(k\) in (F1T38.2). Expanding the actual tensor power of
the direct sum of the two Weil modules, choose a summand
containing \(W\). Its tensor factors have parameters
\(t_1,\ldots,t_k\in\mathbf F_q^\times\). Let \(U\) be
the nondegenerate quadratic space with Gram matrix
\(\operatorname{diag}(t_1,\ldots,t_k)\). The basis map
on delta functions identifies this tensor summand with
the genuine oscillator module
\[
 \mathcal R_U=\mathbf C[\operatorname{Hom}(X,U)],
 \qquad A=X\oplus X^*,\quad\dim X=r,
 \qquad D=\dim\mathcal R_U=q^{rk}.
 \tag{F1T38.7}
\]
It intertwines \(H\) on both sides: on the Weyl-operator
bases this is the direct orthogonal-sum tensor map, and
the residual scalar character of the perfect \(H\) is
trivial. This is also the explicit map in
Montealegre-Mora--Gross, Corollaries 2.2--2.3.
The subgroup \(H\times S\), where \(S=O(U)\) is the
full isometry group, acts through its actual tensor map
to \(\operatorname{Sp}(A\otimes U)\). Any scalar kernel
of that tensor map is retained in its action.

For an isotropic subspace \(N\subset U\) of dimension
\(j\), let \(U_N=N^\perp/N\). The explicit normalized map
\[
 i_N:\mathbf C[\operatorname{Hom}(X,U_N)]
       \longrightarrow\mathcal R_U,
 \qquad
 \delta_{[F]}\longmapsto q^{-rj/2}
                  \sum_{\Delta\in\operatorname{Hom}(X,N)}
                         \delta_{F+\Delta}
 \tag{F1T38.8}
\]
is independent of the lift \(F:X\to N^\perp\), is an
isometry, and is an \(H\)-intertwiner. Its image
\(C_N\) consists exactly of the functions supported on
\(\operatorname{Hom}(X,N^\perp)\) and constant on the
indicated cosets. To check the action, take a Witt
decomposition \(U=\mathbb H^j\perp U_N\). For each
hyperbolic plane the two opposite-parameter Weil factors
identify with the conjugation representation on an
endomorphism space. The normalized identity operator is
fixed; its delta-function expression is the coset sum
in (F1T38.8). Tensoring these \(j\) identity operators
with the \(U_N\) oscillator module proves the claimed
intertwining. Thus
\[
        \dim C_N=q^{r(k-2j)}.
 \tag{F1T38.9}
\]
Diagonalizing the nondegenerate form on \(U_N\) gives an
actual tensor product of \(k-2j\) Weil modules. T37's
square-class intertwiners put it inside the corresponding
power of the original level generator.

\paragraph{3. The stable lower-rank assertion with its exact Fourier maps.}
We give the needed form of the published
Montealegre-Mora--Gross Theorem 1.2, together with its
proof, so its precise use is explicit. The Siegel
unipotent weight of a delta function \(\delta_F\) in
(F1T38.7) is the symmetric form \(\frac12F^*\beta_U F\).
Call the rank of such a weight its ordinary matrix rank.
For an irreducible \(H\)-submodule, its weight rank is the
largest rank of its weights. If this rank is \(d<k\)
and \(r\ge k\), the module is contained in the sum of
the spaces \(C_N\) with \(\dim N=(k-d)/2\); in
particular \(k-d\) is even.

Here are the Fourier and incidence details of that
assertion. For a map \(F:X\to U\) of weight rank \(d\),
put \(N_F=\operatorname{im}F\cap(\operatorname{im}F)^\perp\).
A nondegenerate complement of \(N_F\) in \(\operatorname{im}F\)
has dimension \(d\), so
\[
                  \dim N_F\le\lfloor(k-d)/2\rfloor.
 \tag{F1T38.10}
\]
Let \(K\) be an \(H\)-stable space of weight rank at most
\(d\), and \(\Phi\in K\). If \(F\in\operatorname{supp}\Phi\)
has weight rank \(d\), and \(\Delta:X\to N_F\) satisfies
\(\operatorname{im}(F|_{\ker\Delta})=\operatorname{im}F\),
choose a complement \(X_2\subset\ker\Delta\) to
\(X_1=\ker F\). The partial Fourier transform in \(X_1\)
is an actual oscillator operator from the embedded
symplectic subgroup. Its transformed slice at \(F|_{X_2}\)
has support only in \(\operatorname{Hom}(X_1,N_F^\perp)\):
adding a vector not orthogonal to the radical \(N_F\)
would increase the rank of the form on \(\operatorname{im}F\).
Finite additive-character orthogonality says that Fourier
support in that subspace is equivalent to invariance of
the original slice under \(\operatorname{Hom}(X_1,N_F)\).
Consequently \(\Phi(F+\Delta)=\Phi(F)\).

This invariance holds along every part of a coset where
the radical remains a fixed isotropic \(N\). To see the
extension explicitly, take \(F,F'\) with radical \(N\),
weight rank \(d\), and \(F-F'\in\operatorname{Hom}(X,N)\).
Choose \(\Delta\) uniformly in \(\operatorname{Hom}(X,N)\).
Its kernel has dimension at least
\(r-\dim N\ge d+\dim N=\operatorname{rank}F\).
Choose in it a subspace of dimension \(z=\operatorname{rank}F\).
By invariance under \(GL(X)\), this subspace is uniformly
distributed among all \(z\)-subspaces of \(X\).
The probability of meeting \(\ker F\) nontrivially is at
most
\[
 \frac{q^{r-z}-1}{q-1}\frac{q^z-1}{q^r-1}
                         <\frac1{q-1}.
 \tag{F1T38.11}
\]
The same estimate applies to \(F'\) and
\(\Delta+F-F'\). Since \(q\ge3\), the union bound gives
a choice for which both local invariances apply.
The intermediate map \(F+\Delta\) then has the same
value of \(\Phi\) as both \(F\) and \(F'\).

Now choose a nonzero vector \(\Phi\) of a maximal-rank
weight in an irreducible module of rank \(d\). If
\(k-d\) is even, subtract its coset sums on every branch
with \(\dim N_F=(k-d)/2\), using the just-proved
invariance to define each coefficient. Distinct such
branches do not interfere: a map of this radical
dimension with image in \(N^\perp\) must have image
exactly \(N^\perp\), and hence radical \(N\).
The remainder \(\Phi'\) has the same weight as \(\Phi\).
It belongs to \(K+\sum_N C_N\), an \(H\)-stable space
of rank at most \(d\). If \(k-d\) is odd, make no
subtraction and set \(\Phi'=\Phi\). In both cases, every
map \(F\) in its support satisfies
\[
                         k-2\dim N_F>d.
 \tag{F1T38.12}
\]
Suppose the remainder is nonzero and choose \(F\) in
its support of maximal ordinary map rank. Set
\(X_1=\ker F\) and take any complement \(X_2\).
The slice \(G\mapsto\Phi'(F+G)\), with
\(G:X_1\to U\), is supported on
\(\operatorname{Hom}(X_1,N_F)\). Indeed, maximality
forces \(\operatorname{im}G\subset\operatorname{im}F\),
while equality of the weights forces it to be
orthogonal to \(\operatorname{im}F\). The local invariance
above makes this slice a nonzero constant times the
indicator function of that subspace. Its partial Fourier
transform is therefore a nonzero constant times the
indicator function of \(\operatorname{Hom}(X_1,N_F^\perp)\).
By (F1T38.12), the quadratic space \(N_F^\perp/N_F\)
has rank greater than \(d\). It contains the nondegenerate
rank-\(d\) part of \(\operatorname{im}F\), and has a
nonzero nondegenerate orthogonal complement. An
anisotropic vector in that complement, used as the image
of a rank-one map \(X_1\to N_F^\perp\), raises the
weight rank to \(d+1\). Here \(X_1\ne0\), because
\(\operatorname{rank}F\le(k+d)/2<k\le r\).
This contradicts the rank bound on the stable space
containing the remainder. Thus \(\Phi'=0\).
The orbit of \(\Phi\) spans the irreducible module and
every \(C_N\) is \(H\)-stable, proving the assertion.

It follows in particular that the whole lower-rank
submodule \(\mathcal R_U^{<k}\) is the sum of the codes
for isotropic lines. Every higher-dimensional isotropic
\(N\) contains a line \(\ell\), and the literal
support and invariance conditions give \(C_N\subset C_\ell\).
Conversely each line code has rank at most \(k-2\).
Hence
\[
 \dim\mathcal R_U^{<k}
 \le\frac{q^k-1}{q-1}\,q^{r(k-2)}
 <E:=q^{k+r(k-2)}.
 \tag{F1T38.13}
\]
If \(W\) had weight rank less than \(k\), its inclusion
in this sum and semisimplicity would give a nonzero map
from a line code onto \(W\). The adjoint for invariant
Hermitian metrics is an injection of \(W\) into that
code. Formula (F1T38.8) would then place \(W\) in a
tensor power of level \(k-2\), contrary to minimality.
Thus its weight rank is exactly \(k\).

\paragraph{4. A top constituent of multiplicity one and the actual projector.}
Write the actual \(S\)-decomposition as
\[
 \mathcal R_U=\bigoplus_{\tau\in\operatorname{Irr}(S)}
                 E_\tau\otimes D_\tau,
 \qquad D_\tau=\operatorname{Hom}_S(E_\tau,\mathcal R_U),
 \qquad a_\tau=\dim E_\tau.
 \tag{F1T38.14}
\]
Gurevich--Howe Theorem 4.3.3 gives a unique rank-\(k\)
irreducible constituent \(\eta_\tau\) in each \(D_\tau\),
with multiplicity one; all other constituents have
rank less than \(k\). The following recalls the complete
map establishing that multiplicity. A rank-\(k\) weight
is represented by a surjection \(F:X\to U\).
Its weight space has delta basis the maps \(F'\) with
\(F'^*\beta_UF'=F^*\beta_UF\). The unique factorization
through \(X/\ker F\) makes these maps an \(S\)-torsor.
The stabilizer of this form in \(GL(X)\) acts on the
torsor through the other regular action of
\(O(X/\ker F)\), and maps onto that group. The two
regular actions generate mutual commutant algebras.
The scalar determinant characters in the actual
Schrodinger operators multiply these regular operators
by nonzero scalars and leave the generated algebras
unchanged. Therefore the weight space of \(D_\tau\)
is an irreducible module of dimension \(a_\tau\) for
this stabilizer. Its decomposition induced from any
irreducible \(H\)-decomposition of \(D_\tau\) has exactly
one nonzero summand. All rank-\(k\) weights are in the
same \(GL(X)\)-orbit of this quadratic type. This proves
the asserted uniqueness and multiplicity one.

Choose \(\tau\) such that \(W\simeq\eta_\tau\), put
\(a=a_\tau\), and write \(D_\tau=\eta_\tau\oplus B_\tau\).
By (F1T38.13), \(\dim B_\tau<E\). With \(s=|S|\), the
actual central projector and its trace are
\[
 p_\tau=\frac a s\sum_{u\in S}\overline{\tau(u)}\,\mathcal R_U(u),
 \qquad
 \operatorname{Tr}D_\tau(x)
   =\frac1s\sum_{u\in S}\overline{\tau(u)}\,
                            \operatorname{Tr}\mathcal R_U(x\otimes u).
 \tag{F1T38.15}
\]
The second identity follows by dividing
\(\operatorname{Tr}(p_\tau\mathcal R_U(x))\) by the
actual multiplicity \(a\); it is not the trace of the
whole \(S\)-isotypic summand. This fixes the normalization
needed for the error estimate.

\paragraph{5. Uniform traces at the rank-two commutator.}
The exact Weyl-operator basis of
\(\operatorname{End}(\mathcal R_U)\) is indexed by
\(A\otimes U\). Conjugation by an oscillator operator
permutes this basis by the corresponding symplectic
map. As in T37, taking its trace gives
\[
 \bigl|\operatorname{Tr}\mathcal R_U(z)\bigr|^2
             =q^{\dim\ker(z-I)}
 \quad\bigl(z\in\operatorname{Sp}(A\otimes U)\bigr).
 \tag{F1T38.16}
\]
For \(u\ne I\) in \(S\), the fixed space of \(I_A\otimes u\)
has dimension at most \(2r(k-1)\). Moreover the actual
subspace \(\ker(h-I)\otimes U\) is stable under
\(h\otimes u\), and the restriction of \(h\otimes u-I\)
to it has rank
\((2r-2)\operatorname{rank}(u-I)\ge2r-2\).
At \(u=I\), (F1T38.6) gives rank \(2k\). Thus
\[
 \begin{split}
 |\operatorname{Tr}\mathcal R_U(I\otimes u)|
                 &\le q^{r(k-1)}\quad(u\ne I),\\
 |\operatorname{Tr}\mathcal R_U(h\otimes I)|
                 &=q^{rk-k},\\
 |\operatorname{Tr}\mathcal R_U(h\otimes u)|
                 &\le q^{rk-r+1}\quad(u\ne I).
 \end{split}
 \tag{F1T38.17}
\]
These statements include nonsemisimple elements and use
their actual kernels, rather than only their eigenvalues.

Since \(|\tau(u)|\le a\) and
\(s\le|GL_k(q)|<q^{k^2}\), set \(B=aD/s\). Equations
(F1T38.13)--(F1T38.17) give
\[
 \begin{split}
 \dim\eta_\tau
  &>B\bigl(1-q^{k^2-r}-q^{k^2+k-2r}\bigr),\\
 |\operatorname{Tr}\eta_\tau(h)|
  &<B\bigl(q^{-k}+q^{k^2-r+1}+q^{k^2+k-2r}\bigr).
 \end{split}
 \tag{F1T38.18}
\]
For the error term, explicitly,
\(E/B=(s/a)q^{k-2r}\le q^{k^2+k-2r}\);
using the whole lower-rank dimension only weakens the
estimate. The displayed denominator is positive in the
range below.

Our actual rank satisfies \(r\ge L_0^2+5\ge k^2+5\).
For every \(q\ge3\) and integer \(k\ge2\), the quotient
of the two brackets in (F1T38.18) is at most
\[
 \frac{1/9+1/81+3^{-12}}{1-1/243-3^{-12}}
                         <\frac18.
 \tag{F1T38.19}
\]
Indeed, the three error exponents are at most
\(-5,-4,-12\), respectively. After multiplying by the
positive denominator, the last strict inequality has
the exact positive gap
\[
 1-1/243-3^{-12}-8(1/9+1/81+3^{-12})
                         =\frac{485}{59049}>0.
 \tag{F1T38.20}
\]
This proves the estimate on every original irreducible
constituent, without an assertion that an outer action
stabilizes its theta label.

\paragraph{6. Original multiplicities, phases and the branch contradiction.}
The averaged central \(H\)-projectors and the actual
evaluation maps give
\(V_{\mathbf C}=\bigoplus_i W_i\otimes
\operatorname{Hom}_H(W_i,V_{\mathbf C})\), as in T36.
The particular element \(h\in H\) in (F1T38.6) acts
on each multiplicity space as the identity. Applying
(F1T38.19) to every \(W_i\) and summing with these
actual multiplicities gives
\[
       |\operatorname{Tr}_{V_{\mathbf C}}\beta(h)|<m/8.
 \tag{F1T38.21}
\]
Choose a \(G\)-invariant positive Hermitian metric.
For a unitary operator, each nonfixed eigenvalue
contributes at most two to \(m-\operatorname{Re}\operatorname{Tr}\),
so
\[
 \operatorname{rank}(\beta(h)-I)>7m/16.
 \tag{F1T38.22}
\]
For \(e=\beta(t_{v,c})\), the actual identity
\[
 [g,e]-I=(g-I)e g^{-1}e^{-1}+e(g^{-1}-I)e^{-1}
 \tag{F1T38.23}
\]
bounds this same rank by \(2\operatorname{rank}(g-I)\).
Thus every nonscalar original \(g\) has moving rank
strictly greater than \(7m/32\). A nonidentity scalar
has rank \(m\) and satisfies the same inequality.
The original nonidentity inertia tuple, with \(b\ge12\),
therefore gives
\[
 2m+6=\sum_{j=1}^b\operatorname{rank}(g_j-I)
                  >\frac{7bm}{32}\ge\frac{21m}{8}.
 \tag{F1T38.24}
\]
It follows that \(m<48/5\), contrary to \(m\ge27\).
This proves (F1T38.1).

\paragraph{7. Source, corner and framework maps.}
The group elements in (F1T38.4)--(F1T38.6) map through
the specified \(\beta\) to commutators of the original
monodromy matrices. The tensor and code maps are
auxiliary embeddings used to estimate those matrices;
the original coefficient representation is never
replaced by a tensor factor or by the whole oscillator
module. Its actual constituent multiplicities are
retained in (F1T38.21). The original curve cover,
parabolic cochain complex, six-dimensional primitive
corner and cup-product Hermitian pairing are carried
by identity maps, with polarization multiplier one.
The same identities preserve its integral lattice,
Poincare correspondence and projector, Hilbert/Chow
cycle maps, variational cycle classes and Adams--GRR
maps recorded in F01--F04. No detector component is
added. This restriction concerns precisely the odd
symplectic almost simple branch; the even symplectic,
orthogonal and other stated source gates retain their
existing scopes.

\paragraph{Primary sources and supplied-proof scope.}
The actual cover and outer maps are the retained T37
maps from Malle--Testerman and Steinberg, and the level
bound is the checked T36 application of
Guralnick--Larsen--Tiep, \emph{Inventiones mathematicae}
\textbf{235} (2024), Definition 3.2 and Theorem 3.5.
The top-component construction is Gurevich--Howe,
\emph{Representation Theory, Number Theory, and
Invariant Theory}, Progress in Mathematics 323 (2017),
Theorem 4.3.3 and its complete Appendix 2 proof,
printed pp.~222--224, 229--232.
The code construction and stable lower-rank theorem
are Montealegre-Mora--Gross, \emph{Representation Theory}
\textbf{25} (2021), 193--223, Corollaries 2.2--2.3,
Lemma 2.7, Theorem 1.2 and its complete Section 3.2
proof, with the Fourier calculation in Appendix A.2.
Those operative proofs were read in full in the
published versions. Paragraphs 2--4 give the needed
maps and a direct reconstruction, including the
precise support/invariance Fourier identity and the
stable space containing the subtracted vector.
The separate source record states the inspected
comparison scope of later theta-correspondence papers;
none of their additional exhaustion assertions is a
premise here.
